\documentclass{article} % For LaTeX2e


%%%
% STANDARD PREAMBLE
%%%
%https://tex.stackexchange.com/questions/68821/is-it-possible-to-create-a-latex-preamble-header
\usepackage{/Users/miw267/Repos/latex_preamble/preamble}


%%%
% SPECIAL TO THIS DOC
%%%

\usepackage{amsthm}
\usepackage{tcolorbox}
\tcbuselibrary{theorems}

% Colored theorems
\let\theorem\relax
\let\endtheorem\relax
\newtcbtheorem[number within=section]{theorem}{Theorem}
{colback=blue!5,
 colframe=blue!70!black,
 fonttitle=\bfseries}{thm}
 
 %%%
 % START DOC
 %%%
 
\title{Diffusion Processes: Notes}
%\author{Mike Wojnowicz}


\begin{document}

\maketitle

\begin{abstract}
These notes largely follow \citet[Chapter 6]{chang2007stochastic}.

\end{abstract}

\section{Kolmogorov Backward Equation}

\subsection{Claim} \label{sec:claim}

Let $X_t$ be a one-dimensional diffusion process with infinitesimal mean (a.k.a drift coefficient) $\mu(X_t)$ and infinitesimal variance (a.k.a diffusion coefficient) $\sigma^2(X_t)$. Let $V(X_t)$ be some test function (which we interpret as the \qq{value}) of the state at time $t$.   Define the expected value at some terminal time $T$ given   current state $x$ at current time $t$  by:
\[ u(t,x) \defeq \E[V(X_T) \mid X_t = x]. \]
Then this function $u(t,x)$ satisfies the \textbf{Kolmogorov Backward Equation}:
 \[ 0 = u_t + u_x \mu(x) + \half u_{xx} \sigma^2(x).  \]
The boundary condition is given by $u(T,x)=V(x)$.
\subsection{Remarks}



Why is this claim interesting?  

\begin{enumerate}
	\item \textit{It is the foundation of option pricing.} In the case of geometric Brownian motion, the backward equation becomes the Black–Scholes PDE.  That is, pricing a derivative is done by solving a backward equation.	
\item \textit{We can use it to find the probability density at any time.}  As a corollary to the claim, the same backward equation holds when $u(t,x)$ is interpreted in terms of the transition density.  Using this, we can find probability density $p(t,x)$ at every time $t>0$.	This highlights an important property of diffusion processes:

\fbox{%
\begin{minipage}{0.9\linewidth}
\begin{center}
	The drift $\mu(x)$ and diffusion $\sigma^2(x)$ coefficients uniquely determine the probability density $p(t,x)$ at every time $t>0$.
\end{center}
\end{minipage}%
}
	




%As a corollary, the same backward equation holds when $u(t,x)$ is interpreted in terms of the transition density.

%\[ f(t',y \mid t,x) = \frac{\partial}{\partial y} P\bigg(X(t') \leq y \mid X(t)=x \bigg) \]
 %So we only need the drift and diffusion coefficients $\big(\mu(x),\sigma^2\big(x))$ to determine the density $f$.  Moreover, we solve this deterministically (via a PDE) rather than requiring any stochastic simulation.
\end{enumerate}


Why is this called the \qq{backward equation}?

\begin{enumerate}
\item The PDE runs backward in time from $T$ to earlier times $t$.   (Note that it differentiates with respect to the \textit{starting} time $t$.)
\end{enumerate}


\subsection{Heuristic Justification of Claim~\ref{sec:claim}.}

\textbf{Step 1. Consider discretized diffusion}
 

Let $\Delta X = X_{t + \Delta t} - X_t$ be a small change in the diffusion process. Then per \citet[pp.180]{chang2007stochastic}, we have
%
\begin{align}
 \Delta X = \mu(X_t) \Delta t + \sigma(X_t) \, \sqrt{\Delta t} \, Z  \qquad Z, \sim N(0,1)
 \label{eqn:discretized_diffusion}
\end{align}

Note in particular that, as given in \citet[Eqn. 6.2 and Exercise 6.4]{chang2007stochastic}, we have
\begin{subequations}
\begin{align}
\E[\Delta X \mid X(t)=x] = \mu(x) \Delta t + o(\Delta t) \label{eqn:expected_delta_x}\\	
\E[(\Delta X)^2 \mid X(t)=x] = \sigma^2(x) \Delta t + o(\Delta t) \label{eqn:expected_delta_x_squared} 
\end{align}
as $\Delta t \downarrow 0$.  (Exercise: defend these equations.)
\end{subequations}




\textbf{Step 2. Use the Law of Iterated Expectation}
 
We have 
%
\begin{align*}
u(t,x) &= \E[V(X_T) \mid X_t=x] && \scripttext{(definition)} \\
&= \E\big[\E[V(X_T) \mid X_t=x, X_{t+\Delta t}] \mid X_t = x\big] && \scripttext{(LIE)} \\
&= \E\big[\E[V(X_T) \mid \cancel{X_t=x}, X_{t+\Delta t}] \mid X_t = x\big] && \scripttext{(Markov property)} \\
&= \E \big[u( t+ \Delta t, X_{t+\Delta t}) \mid X_t = x\big] && \scripttext{(definition)} \\
\labelit\label{eqn:expected_value}
\end{align*}

\fbox{%
\begin{minipage}{0.9\linewidth}
\begin{center}
Note: This equation is \underline{dynamic programming}, expressing the fact that

\begin{quote}
Starting now is equivalent to starting an instant later and averaging over the random step in between.	
\end{quote}
\end{center}
\end{minipage}

}

\textbf{Step 3. Form 2nd-order Taylor series expansion}

We have 
%
\begin{align*}
u(t+ \Delta t, X_{t+\Delta t}) & = u(t,x)  + u_x(t,x) \Delta X + u_t(t,x) \Delta t + \half u_{xx}(t,x) (\Delta X)^2+ o(\Delta t),  
\end{align*}
%
where $o(\Delta t)$ absorbs the other second-order terms $\half u_{tt}(\Delta t)^2$ and $u_{tx}(\Delta t \Delta X)$. 

\vspace{1cm}

{\footnotesize
\begin{itemize}
\item \underline{Details:} Here are some details on why $o(\Delta t)$ absorbs the second-order term $u_{tx}(\Delta t \Delta X)$. First note from \Eqref{eqn:discretized_diffusion}  that over a small time interval 
$\Delta t$, the typical size of the change in $X$ is proportional to $\sqrt{\Delta t}$.  In particular, we have
%
\begin{align}
 \Delta X = O_p(\sqrt{\Delta t}) 
\label{eqn:space_increment_scales_quadratically_in_time_increments}	
\end{align}
%
which means that $\frac{\Delta X}{\sqrt{\Delta t}}$ is bounded in probability as $\Delta t \downarrow 0$. This is known as the \qq{quadratic variation} property of diffusion.


Now multiply the orders:

\[ \Delta t \Delta X = \Delta t O_p(\sqrt{\Delta t}) = O_p(\Delta t^{3/2}).\]

Now since
\[ \frac{\Delta t^{3/2}}{\Delta t} = \Delta t^{1/2} \to 0, \]

we may conclude
\[ \Delta t \Delta X =o_p(\Delta t), \]
since multiplying by the bounded coefficient does not change the order.  Note that $o_p$ refers to \qq{little-o in probability}.  In these notes, we write $o_p$ as $o$, as we simply want to give the gist.
\end{itemize}
}

\textbf{Step 4. Substitute}

Now we substitute the Taylor series expansion (Step 3)  into the RHS of \Eqref{eqn:expected_value}:
%
\begin{align*}
\E \big[&u( t+ \Delta t, X_{t+\Delta t}) \mid X_t = x\big] \\
&= \E \big[ u(t,x)  + u_t(t,x) \Delta t + u_x(t,x) \Delta X + \half u_{xx}(t,x) (\Delta X)^2+ o(\Delta t)  \mid X_t = x\big] + o(\Delta t) && \scripttext{(substitution)} \\
&= u(t,x)  + u_t(t,x) \Delta t  + u_x(t,x) \E[\Delta X \mid X_t = x] + \half u_{xx}(t,x) \E[(\Delta X)^2 \mid X_t = x] + o(\Delta t) && \scripttext{(linearity)}  \\
&= u(t,x)  + u_x(t,x)  + u_t(t,x) \Delta t + \mu(x) \Delta t +  \half u_{xx}(t,x) \sigma^2(x)\Delta t + o(\Delta t) && \scripttext{(Eqs.\ref{eqn:expected_delta_x},\ref{eqn:expected_delta_x_squared})} \\
&= u(t,x)  + \bigg(u_t(t,x) + u_x(t,x) \mu(x) + \half u_{xx}(t,x) \sigma^2(x) \bigg) \Delta t+  o(\Delta t) && \scripttext{(associative prop.)}  
\end{align*}
%
Now using this within  \Eqref{eqn:expected_value}, we subtract $u(t,x)$ from both sides, divide both sides by $\Delta t$, and take the limit as $\Delta t$ goes to 0 to obtain the result.



\bibliography{references}{}
\bibliographystyle{apalike}

\end{document}


